\chapter{\IfLanguageName{dutch}{Setup}{Setup}}%
\label{ch:setup}

\section{{Docker starten vanuit C\#}}

Binnen het EMAV-project wordt een gestandaardiseerde opzet gehanteerd voor het 
initialiseren van de applicatieomgeving.
Een eerste stap in dit proces is het kopiëren van de noodzakelijke dynamic-link 
libraries (DLL's).
Dankzij een doordachte projectstructuur en consistente naamgevingsconventies 
kunnen deze DLL-bestanden,
die geautomatiseerde manier gegenereerd worden tijdens het buildproces 
van de applicatie,
op een geautomatiseerde manier naar de correcte directory worden verplaatst.

De jupyter notebooks worden gecontaineriseerd met behulp van een Dockerfile.
Deze Dockerfile is gebaseerd op een runtime image van Microsoft, 
meer bepaald de .NET 10.0 runtime.
Deze image steunt op een Debian-gebaseerde Linuxdistributie, 
wat een stabiele en breed ondersteunde basis biedt.

Tijdens de buildfase van de container worden bijkomende afhankelijkheden 
geïnstalleerd via het pakketbeheersysteem \texttt{apt-get},
waaronder Python en gerelateerde packages.
Aansluitend worden de eerder gegenereerde DLL-bestanden in de container 
gekopieerd, zodat deze beschikbaar zijn voor verdere verwerking.

Na het opzetten van de omgeving wordt een Jupyter-server opgestart via een 
specifiek startupscript.
Dit script zorgt ervoor dat de relevante DLL's worden ingeladen 
en dat de belangerijkste namespaces toegankelijk worden gemaakt binnen de Python-omgeving.
Voor het gebruik van specifieke modellen of functionaliteiten kunnen bijkomende import vereist zijn.

Deze aanpak maakt het mogelijk om een reproduceerbare en consistente ontwikkel- 
en uitvoeringsomgeving te garanderen, wat essentieel is binnen een 
professionele software ontwikkelingscontext.


Voor de interop tussen C\# en docker wordt er gebruik gemaakt van Docker.Dotnet \textcite{docker_dotnet}.

Voor het automatisch opzetten van een Jupyter-gebaseerde 
analyseomgeving werd een \texttt{JupyterSetup}-klasse geimplementeerd in C\#.
Deze klasse is verantwoordelijk voor twee hoofdtaken:
\begin{itemize}
	\item het verzamelen en kopieren van de benodigde DLL-bestanden en
	\item het opstarten van een Docker-container waarin een Jupyter-server draait
\end{itemize}

\subsection{Algemene opzet}
De centrale methode van deze klasse is de \texttt{Setup}-methode, die beide stappen sequentieel uitvoert. De implementatie hiervan is terug te vinden in Codefragment \ref{lst:setup_main}.

\begin{listing}[ht]
\begin{minted}{csharp}
public static async Task Setup()
{
    SetupDLL();
    await SetupDocker();
}
\end{minted}
\caption{De centrale \texttt{Setup}-methode van de \texttt{JupyterSetup}-klasse.}
\label{lst:setup_main}
\end{listing}

Hierbij wordt eerst de lokale omgeving voorbereid (DLL's), waarna de containeromgeving wordt opgestart.

\subsection{Kopieren van DLL-bestanden}
De methode \texttt{SetupDLL} is verantwoordelijk voor het verzamelen van alle relevante DLL-bestanden die tijdens het buildproces gegenereerd worden.
Op basis van de locatie van de uitvoerende assembly wordt de root van de solution bepaald. De volledige logica voor het filteren en kopiëren van deze bestanden is opgenomen in Bijlage \ref{ch:codefragmenten} (zie Codefragment \ref{lst:setupdll}).

De aanpak is eenvoudig maar robuust: alle DLL's uit de \texttt{obj}-directories 
(build-output) worden verzameld,
maar referentie-assemblies en runtime-specifieke bestanden worden uitgesloten.
Dit voorkómt duplicaten en onnodig grote container-images. 
De gekopieerde bestanden gaan naar \texttt{ILVO.EMAV.Jupyter/dlls},
die later in de Docker-image worden opgenomen.

\section{Docker-orchestratie via Docker.DotNet}

\subsection{Verbinding met Docker}

De \texttt{Docker.DotNet}-bibliotheek abstraheert de Docker-daemon-communicatie. 
Op Windows gebruikt deze named pipes,
op Linux Unix-sockets. De klasse initialiseert eenmalig de client, zoals getoond in Codefragment \ref{lst:docker_init}:

\begin{listing}[ht]
\begin{minted}{csharp}
docker = new DockerClientConfiguration(
    new Uri(RuntimeInformation.IsOSPlatform(OSPlatform.Windows)
        ? "npipe://./pipe/docker_engine"
        : "unix:///var/run/docker.sock")
).CreateClient();
\end{minted}
\caption{Initialisatie van de \texttt{DockerClient} voor verschillende platformen.}
\label{lst:docker_init}
\end{listing}

\subsection{Container aanmaken en poorten binden}

De container wordt aangemaakt op basis van de \texttt{jupyter-image}. Via dynamische poortbinding
kunnen meerdere Jupyter-instanties parallel op dezelfde machine draaien. De gedetailleerde configuratie van de \texttt{CreateContainerParameters} is te vinden in Codefragment \ref{lst:docker_create}.

Wanneer \texttt{HostPort} leeg is, kiest Docker een vrije poort automatisch. 
Dit voorkómt conflicten en maakt
het mogelijk dat onderzoeksgroepen parallel hun eigen analyseomgeving draaien.

\subsection{Token- en poort-extractie uit logs}

Jupyter genereert bij opstarten een beveiligingstoken in de logs. 
Dit moet aan gebruikers ter beschikking gesteld worden.
De implementatie parseert de container-logs met reguliere expressies, zoals weergegeven in Codefragment \ref{lst:docker_logs}.

De \texttt{DemultiplexStream}-methode zorgt ervoor dat de Docker-logs 
(die in een gedemultiplexed formaat komen)
correct naar UTF-8 tekst worden omgezet. Deze URL kan vervolgens aan de gebruiker gegeven worden.

\section{Jupyter-container configuratie}

\subsection{Docker base-image en Python-setup}

De Dockerfile bouwt voort op de .NET 10.0 runtime, wat native .NET-ondersteuning in Python garandeert. Het volledige build-script is opgenomen in Bijlage \ref{app:container_config} (zie Codefragment \ref{lst:dockerfile}).

De \texttt{--break-system-packages}-vlag is nodig omdat deze container geen systeem-wide Python is,
en pip's veiligheidsmechanismen dus niet relevant zijn.

\subsection{Automatische initialisatie van .NET-referenties}

IPython voert automatisch configuratie uit wanneer het opstart. Dit wordt gebruikt om alle .NET-DLL's
in te laden zonder dat gebruikers dat handmatig moeten doen. De specifieke configuratie hiervoor is te vinden in Codefragment \ref{lst:startuppy}.

Hierdoor zijn alle klassen en functies uit de domein- en infrastructuurlaag onmiddellijk beschikbaar
in notebooks zonder extra import-statements.

\section{Veiligheids- en isolatiemaatregelen}

\subsection{Netwerkisolatie}

De Jupyter-containers draaien in een afgesloten Docker-netwerk. Ze kunnen enkel communicate met de API-laag
(de ASP.NET Razor Pages applicatie). Dit voorkómt accidentele of opzettelijke verbindingen naar externe systemen.

\subsection{Processesolatie en resourcebeperkingen}

Containers draaien onder een niet-root gebruiker om privilege-escalation aan te remmen.
Resourcebeperkingen (geheugen, CPU, aantal processen) kunnen ingesteld worden om DOS-aanvallen via runaway-notebooks
te mitigeren. Het bestandssysteem is zoveel mogelijk read-only, met schrijfrechten beperkt tot werkdirectories.

\subsection{Vereenvoudigde startprocedure}

De gehele setup wordt aangeboden via een enkele \texttt{Setup()}-methode. Dit zorgt voor eenvoud voor eindgebruikers:
een enkele aanroep initialiseert DLL's en start de container. Dit minimaliseert fouten en handmatige tussenkomst.
